\section{The primitive metric addition}
\label{sec:primitive-statement}

The global construction will add finitely many positive tensors of rank one.
We first state the local step, including the boundary conditions that allow
one step to be followed by the next.  Only a preferred unit normal is
propagated; a full normal frame is chosen locally when needed.

Let $F:M\to\R^4$ be a smooth immersion, and write $\gamma=\gamma(F)$.
Suppose $C_0,\ldots,C_m$ are smooth embedded closed curves with nonzero
smooth vector fields $Y_i$ defined on neighborhoods of $C_i$.
The fields $Y_i$ are not required to be tangent to the curves; in the
application they are the transverse coordinate directions of later phases.
Set
\[
 A_i(F)=B_F(Y_i,Y_i)\quad\hbox{on }C_i.
\]
For $i\ne l$, wherever the curves intersect and $A_i(F)\ne0$, define
\begin{equation}
 Q_{il}(F)=
 \left\langle B_F(Y_i,Y_l),\frac{A_i(F)}{|A_i(F)|}\right\rangle^2
 +K_{\gamma}|Y_i\wedge Y_l|_{\gamma}^2.
 \label{eq:crossing-invariant}
\end{equation}
Here $K_\gamma$ is Gauss curvature, and
$|u\wedge v|_\gamma^2=\gamma(u,u)\gamma(v,v)-\gamma(u,v)^2$.
The Gauss equation gives a useful interpretation.  With
$\widehat A_i=A_i(F)/|A_i(F)|$ and $Z_{il}=B_F(Y_i,Y_l)$,
\begin{equation}
 Q_{il}(F)=\langle A_i(F),A_l(F)\rangle
       -\bigl|Z_{il}-\langle Z_{il},\widehat A_i\rangle\widehat A_i\bigr|^2.
 \label{eq:crossing-gauss-interpretation}
\end{equation}
Thus $Q_{il}>0$ forces the two pure second-form vectors to have positive
inner product.  It also shows that $Q_{il}$ depends continuously on the
$2$-jet of $F$ wherever $A_i\ne0$, despite its curvature expression
in \eqref{eq:crossing-invariant}.  Both observations will be used when
the preferred normal is joined across the support boundary.
For a preferred normal $n$, the conditions to be preserved are
\begin{equation}
 \begin{aligned}
 &A_i(F)\ne0,\qquad A_i(F)/|A_i(F)|\ne-n
 &&\text{on }C_i,\\
 &\langle A_i(F),n\rangle>0,\qquad Q_{il}(F)>0
 &&\text{on }C_i\cap C_l\quad(i\ne l).
 \end{aligned}
 \label{eq:boundary-invariants}
\end{equation}
The crossing conditions are imposed in both orders.  In particular they
require positivity of the dot product with $n$ for each of the two pure
second-form vectors.  The numbers $Q_{il}$ do not depend on the chosen
preferred normal.

\begin{proposition}[Primitive addition with preserved boundary conditions]
\label{prop:primitive-addition}
Let $M$ be a closed smooth surface and $F:M\to\R^4$ a smooth immersion
whose second fundamental form is nonzero as a tensor at every point.
Assume $F$ has a preferred normal $n$.  Let $D$ be a disk with smooth
boundary $C_0$ and coordinates $(x,y)$ on a neighborhood of $\overline D$.
Let $a\in C^\infty(M)$ be nonnegative, positive exactly on $D$, and put
$Y_0=\partial_y$.  Suppose
\begin{equation}
 \Hess_{\gamma(F)}x(\partial_y,\partial_y)>0
 \quad\text{on }\overline D.
 \label{eq:primitive-hessian}
\end{equation}
Let $C_1,\ldots,C_m$ be finitely many later curves with nonzero smooth
vector fields $Y_i$ on their neighborhoods.  Assume the following:
\begin{enumerate}
\item All curves $C_0,\ldots,C_m$ are pairwise transverse and have no
triple intersection.
\item On the portion of each later curve $C_i$, $1\le i\le m$, in $\overline D$, the vector
$\partial_y$ is tangent to $C_i$ at only finitely many points.
\item At every $C_i\cap C_l\cap D$ with distinct $i,l>0$,
$dx(Y_i)\ne0$ and $dx(Y_l)\ne0$.
\item Conditions \eqref{eq:boundary-invariants} hold for all indices
$0,\ldots,m$.
\end{enumerate}
Then, for every $\varepsilon>0$, there are a smooth immersion
$F_+:M\to\R^4$ and a preferred normal $n_+$ such that
\[
 \gamma(F_+)=\gamma_+:=\gamma(F)+a^2\,dx^2,
 \qquad B_{F_+}\ne0\ \text{at every point},
\]
conditions \eqref{eq:boundary-invariants} hold for $(F_+,n_+)$ and all
later indices $1,\ldots,m$, and
\[
 \|F_+-F\|_{C^2(M\setminus D)}<\varepsilon.
\]
The tensor $a^2dx^2$ is extended by zero outside the coordinate neighborhood.
\end{proposition}

Because $a$ is smooth and vanishes outside $D$, it and all its derivatives
vanish on $C_0$.  This makes the extended tensor smooth.  The proposition
allows the map itself to change on the exterior; the stated $C^2$ control
is what will preserve the boundary data of the later steps.

The hypotheses have different roles.  Table~\ref{tab:invariants} records
how they are established initially and passed to the next primitive step.
In particular, the phase-Hessian condition depends on the current metric,
whereas the boundary conditions also depend on the immersion and normal.

\begin{table}[htbp]
\centering
\small
\renewcommand{\arraystretch}{1.18}
\begin{tabular}{@{}>{\raggedright\arraybackslash}p{.21\linewidth}>{\raggedright\arraybackslash}p{.24\linewidth}>{\raggedright\arraybackslash}p{.23\linewidth}>{\raggedright\arraybackslash}p{.23\linewidth}@{}}
\toprule
Condition & Initial construction & Use & Preservation \\
\midrule
$B_F\ne0$ everywhere
& Radial component of a spherical immersion
& Good phases and exact correction
& Nonzero transverse second form on the disk; exterior $C^2$ control \\
Global unit normal $n$
& Negative radial section
& Choice of velocity-plane frame and boundary directions
& Local adjustments, projection and normalized interpolation \\
Positive transverse phase Hessian
& Convex phases and sufficiently many finite cycles
& Independence of $F_y,F_{yy}$
& Each actual partial metric is prescribed in advance \\
$A_i\ne0$ and $A_i/|A_i|\ne-n$
& Positive radial projection
& Velocity frame near each later support boundary
& Gauss identity and the turn construction \\
Positive normal projections and both ordered $Q_{il}>0$ at crossings
& Second form in $S^3$ made zero at the finite crossing set
& Collar estimates and successive boundary control
& Mixed--pure comparison and local normal adjustments \\
\bottomrule
\end{tabular}
\caption{The geometric data carried through the finite construction.
Section~\ref{sec:global-construction} supplies the initial data and verifies
the phase-Hessian condition for every prescribed partial metric.
Sections~\ref{sec:velocity-loop}--\ref{sec:geometry-preservation}
preserve the immersion and normal conditions.  A single unit normal suffices.}
\label{tab:invariants}
\end{table}

We prove the proposition in three parts.  After the analytic correction
results, Section~\ref{sec:velocity-loop} constructs a periodic velocity
whose averaged displacement is zero and whose leading metric adds
$a^2dx^2$.  Section~\ref{sec:primitive-realization} realizes this velocity
by exact smooth immersions with controlled derivatives.
Section~\ref{sec:geometry-preservation} proves the boundary conditions and
constructs the new global preferred normal.
